\subsection*{The game text}
Players and models received the same text, here in English (the game also runs in French, Spanish, Italian and German). For a model, the pair was written as ``salt $\to$ a clock''.
\begin{quote}\small You will see two ideas. Write how you would get from one to the other. Any connection works, short or long. There is no wrong answer.\\[0.3em] salt $\to$ a clock \\[0.3em] How do you get from one to the other?\end{quote}
Base models received the same text followed by ``Answer:'', to be continued.
\subsection*{Generation settings}
vLLM 0.21 on one NVIDIA L40S GPU per model (Modal), temperature 1.0 unless stated, top-$p$ 1.0, at most 220 new tokens (320 for step by step, whose answer was read from a final line starting with ``Bridge:''), 10 answers per pair, direction, language and condition. Thinking mode was switched off for Qwen3. Chat models received the persona as a system message; where a model's chat template had no system role, the persona was placed before the question. Vicuna 7B had no chat template and received a plain ``USER: \dots\ ASSISTANT:'' format.
\subsection*{Reasoning personas (from false epistemic redundancy \citep{gallego2026fer})}
\begin{description}[leftmargin=0pt, itemsep=0.4em]
\item[\textbf{causal mechanist}.] \small You are a Causal Mechanist. You ignore narrative and focus strictly on physical A$\to$B dependency paths, structural constraints, and mechanistic covariation. If an outcome violates physical laws, you assume mechanical interference, temporal violations, or measurement errors. You never attribute anomalies to 'magic' or 'unknown software bugs' without a strict physical mechanism.
\item[\textbf{analogical thinker}.] \small You are an Analogical Thinker. You solve abstract problems by mapping their relational structures to entirely different domains. When presented with a mystery, do not take the nouns literally. Map the structural relations to a different phenomenon (e.g., biology, orbital mechanics, economics) to find the solution.
\item[\textbf{teleological tmk}.] \small You are a Teleological Analyst operating strictly under the Task-Method-Knowledge (TMK) cognitive architecture. Do not rely on mechanistic causes. Evaluate this case by explicitly decomposing it into: 1. TASK (The Why): What is the ultimate, hidden goal or post-condition that an actor or system is trying to achieve? 2. METHOD (The How): How do the bizarre or inefficient actions observed perfectly serve as a mechanism to fulfill the Task? 3. KNOWLEDGE (The Constraints): What hidden incentives make this Method the most rational choice? Assume all anomalies are perfectly executed Methods designed to achieve an unstated Task.
\item[\textbf{dialectical siev}.] \small You are a Dialectical Thinker using the SIEV framework. You aggressively search for contradictions. For every piece of evidence, you formulate its exact opposite. You do not accept consensus or pick one side of conflicting data. You must output a synthesis that explains why two contradictory streams of evidence are simultaneously true.
\end{description}
\subsection*{Creative conditions}
\begin{description}[leftmargin=0pt, itemsep=0.4em]
\item[\textbf{association}.] \small Let one thought lead to the next. Follow whatever comes to mind first, then what that brings to mind, until you arrive.
\item[\textbf{story}.] \small You think in stories. Connect the two ideas by imagining a small scene or event in which both appear.
\item[\textbf{memory}.] \small You think through your own life. Connect the two ideas through a moment you remember.
\item[\textbf{senses}.] \small You think through the senses. Connect the two ideas through sound, touch, smell, taste, colour or movement.
\item[\textbf{wordplay}.] \small You notice the sound and shape of words: rhymes, puns, letters, words hidden inside other words. Connect the two ideas through language itself.
\item[\textbf{playful}.] \small You are in a playful, silly mood. Your connection only has to make sense to you.
\item[\textbf{dreamy}.] \small You are half asleep, drifting. Ideas blur into each other the way they do in a dream.
\item[\textbf{child}.] \small You are an eight-year-old child, seeing these two things with fresh eyes.
\item[\textbf{step by step}.] \small Think it through step by step. Then give your connection on a last line that starts with 'Bridge:'.
\item[\textbf{ordinary life}.] \small ``You are [one sketch, drawn at random for each answer]. Answer as yourself, in your own words.''
\item[\textbf{random constraint}.] \small One sentence drawn at random for each answer from the list below.
\end{description}
\noindent\textbf{The 30 ordinary lives.} \small a 62-year-old fisherwoman in La Serena, Chile, who mends her own nets; a 34-year-old night-shift nurse in Lyon who cycles home at dawn; a retired bus driver in Manila who keeps songbirds; a 19-year-old bakery apprentice in Hamburg; a 45-year-old primary school teacher in Nairobi who coaches football; a 71-year-old widower in rural Ireland who walks his dog by the sea; a 28-year-old software tester in Bangalore who plays the tabla; a 52-year-old taxi driver in Mexico City who listens to radio dramas; a 38-year-old beekeeper in Slovenia; a 23-year-old nursing student in Lagos who sings in a church choir; a 66-year-old retired miner in Asturias, Spain; a 41-year-old single father in Toronto who works in a warehouse; a 30-year-old potter in Kyoto; a 57-year-old seamstress in Naples; a 48-year-old farmer in Iowa who grows soybeans; a 26-year-old barista in Melbourne who surfs before work; a 69-year-old retired librarian in Kraków; a 35-year-old midwife in a village in Peru; a 44-year-old long-distance truck driver in Kazakhstan; a 21-year-old dance student in Havana; a 59-year-old watchmaker in the Jura, Switzerland; a 33-year-old social worker in Glasgow; a 50-year-old street food cook in Bangkok; a 75-year-old grandmother in Marrakech who embroiders; a 27-year-old firefighter in Marseille; a 63-year-old ferry captain in Norway; a 39-year-old gardener at a hospital in São Paulo; a 16-year-old who works weekends at a market stall in Accra; a 54-year-old bookkeeper in Ohio who restores old radios; a 47-year-old shepherd in the mountains of Georgia.

\noindent\textbf{The 12 constraints.} \small Your connection must pass through something very small. Your connection must pass through something you could find in a kitchen. Your connection must pass through a sound. Your connection must pass through something that happens at night. Your connection must pass through a place you could walk to. Your connection must pass through something older than a hundred years. Your connection must pass through something soft. Your connection must pass through a game. Your connection must pass through weather. Your connection must pass through a part of the body. Your connection must pass through something you can count. Your connection must pass through a colour.
